% Chapter — RCS thrusters and reaction wheels (FR-1). Follows the
% mandatory FR-29 six-section template: purpose; assumptions and domain of
% validity; derivation; discretization and implementation; validation
% evidence; references. The equation labels eq:actuators:mib,
% eq:actuators:rcscoupling, eq:actuators:wheelclamp,
% eq:actuators:wheelreaction, and eq:actuators:totalmomentum are echoed
% verbatim in comments in cpp/src/models/actuators.cpp (FR-29
% equation-label-to-code traceability); renaming them breaks that trace.
\chapter{RCS Thrusters and Reaction Wheels}
\label{ch:actuators}

\section{Purpose}
\label{sec:actuators:purpose}

This chapter documents the FR-1 attitude actuators: on/off reaction
control thrusters grouped in clusters, with minimum-impulse-bit (MIB)
enforcement and full force-and-torque coupling about the current
composite CG, and reaction wheels with per-wheel momentum states and
exact torque and momentum saturation, including the total-system
angular-momentum bookkeeping that makes conservation testable. Both
actuators are implemented in \texttt{cpp/src/models/actuators.cpp} with
the interface in \texttt{cpp/include/star/models/actuators.hpp}.

\section{Assumptions and domain of validity}
\label{sec:actuators:domain}

Assumptions:
\begin{enumerate}
  \item \textbf{Ideal rectangular RCS pulses.} A firing thruster
    delivers its rated constant thrust along a fixed unit direction at
    a fixed position in the vehicle frame (FR-13, $+X$ forward); valve
    opening and closing transients are collapsed into the MIB rule.
    Thruster propellant consumption is not modeled in v1 (RCS
    propellant budgets are orders below main-engine flow for the FR-13
    vehicle class); the delivered momentum is nevertheless exact, so
    conservation checks remain meaningful.
  \item \textbf{Minimum impulse bit as a hard floor.} A commanded pulse
    whose ideal impulse (thrust times duration) falls below the
    configured MIB delivers \emph{exactly zero} impulse --- the valve
    never effectively opens --- and a pulse at or above the MIB
    delivers exactly thrust times duration. Real thrusters deliver a
    degraded, poorly repeatable impulse near the MIB; the model
    replaces that gray zone with a deterministic threshold, documented
    so control designers treat the MIB as unusable resolution.
  \item \textbf{Rigid momentum-exchange wheels.} Each wheel spins about
    a fixed unit axis in the vehicle frame and stores the relative
    angular momentum $h_w$ along that axis; the wheel's contribution to
    the composite inertia is part of the configured dry-body inertia.
    Motor torque is internal to the vehicle: the reaction on the body
    is exactly the negative of the torque on the wheel, following the
    standard momentum-exchange formulation of Markley and
    Crassidis~\cite{markley2014}. Wheel friction and drag torques are
    not modeled.
  \item \textbf{Saturation semantics (documented choice).} Delivered
    wheel torque is clamped at the configured torque limit, and the
    stored momentum is limited at the configured momentum rail: a step
    that would cross the rail delivers only the torque that lands the
    momentum exactly on the rail, and a wheel starting on the rail
    delivers exactly zero torque for a same-sign command (it cannot
    accept further momentum of that sign) while opposite-sign
    (desaturating) commands act normally.
\end{enumerate}

Domain of validity: unit-norm thruster directions and wheel axes;
non-negative thrust, MIB, pulse duration, torque limit, and momentum
limit; strictly positive step size for the wheel update; initial wheel
momentum within $[-h_{\max}, +h_{\max}]$. Out-of-domain behavior,
matching the code exactly: non-finite inputs, malformed unit vectors,
negative ratings or durations, non-positive step sizes, and an initial
momentum outside the rails throw \texttt{std::domain\_error}.

\section{Derivation}
\label{sec:actuators:derivation}

\subsection{RCS pulse impulse and coupling}

A thruster with rated thrust $F$ and unit direction $\hat{\vv{d}}$
fired for a duration $T$ ideally delivers the impulse magnitude $F T$.
With the configured minimum impulse bit $J_{\min}$:
\begin{equation}
  J \;=\;
  \begin{cases}
    0, & F T < J_{\min}, \\
    F T, & F T \ge J_{\min}.
  \end{cases}
  \label{eq:actuators:mib}
\end{equation}
Integrating the constant force and its moment about the composite CG
$\bar{\vv{c}}$ over the pulse (the geometry is constant during the
pulse at the FR-1 modeling level) gives the linear and angular impulse
delivered to the body,
\begin{equation}
  \boxed{\; \Delta\vv{p} \;=\; J\, \hat{\vv{d}}, \qquad
  \Delta\vv{L} \;=\; (\vv{r}_t - \bar{\vv{c}}) \times \Delta\vv{p}, \;}
  \label{eq:actuators:rcscoupling}
\end{equation}
where $\vv{r}_t$ is the thruster position. Because thrusters are
generally not mounted through the CG, every pulse couples force and
torque; a cluster sums \eqref{eq:actuators:rcscoupling} over its
active thrusters in configuration order. While a thruster is held on
continuously the body sees the constant force $F \hat{\vv{d}}$ and
torque $(\vv{r}_t - \bar{\vv{c}}) \times F \hat{\vv{d}}$, the
integrands of \eqref{eq:actuators:rcscoupling}.

\subsection{Reaction-wheel torque delivery and saturation}

A wheel with unit spin axis $\hat{\vv{a}}$ stores the relative angular
momentum $h_w$ along $\hat{\vv{a}}$ and obeys $\dot{h}_w = \tau$, where
$\tau$ is the motor torque on the wheel. Over one step $\Delta t$ with
commanded torque $\tau_{\mathrm{cmd}}$, the delivered torque and new
momentum are, per the saturation semantics of
Section~\ref{sec:actuators:domain}:
\begin{equation}
  \begin{aligned}
    \tau_c &\;=\; \operatorname{clamp}(\tau_{\mathrm{cmd}},\,
      -\tau_{\max},\, +\tau_{\max}), \\
    (\tau, h_w') &\;=\;
    \begin{cases}
      \bigl( (\,+h_{\max} - h_w)/\Delta t,\; +h_{\max} \bigr),
        & h_w + \tau_c \Delta t > +h_{\max}, \\
      \bigl( (\,-h_{\max} - h_w)/\Delta t,\; -h_{\max} \bigr),
        & h_w + \tau_c \Delta t < -h_{\max}, \\
      \bigl( \tau_c,\; h_w + \tau_c \Delta t \bigr),
        & \text{otherwise}.
    \end{cases}
  \end{aligned}
  \label{eq:actuators:wheelclamp}
\end{equation}
On the rail with a same-sign command the first branch yields
$\tau = 0 / \Delta t = 0$ exactly. The motor torque is internal, so
Newton's third law gives the reaction torque on the body
\begin{equation}
  \boxed{\; \vv{\tau}_{\mathrm{body}} \;=\; -\tau\, \hat{\vv{a}}. \;}
  \label{eq:actuators:wheelreaction}
\end{equation}

\subsection{Total angular momentum bookkeeping}

With body inertia $\mat{I}$ (about the composite CG), body angular
velocity $\vv{\omega}$, and $n$ wheels, the total angular momentum of
the vehicle-plus-wheels system, resolved in the body frame, is
\begin{equation}
  \vv{H} \;=\; \mat{I}\, \vv{\omega}
    + \sum_{i=1}^{n} h_{w,i}\, \hat{\vv{a}}_i,
  \label{eq:actuators:totalmomentum}
\end{equation}
the standard momentum-exchange decomposition~\cite{markley2014} (the
wheels' transverse and structural inertia are inside $\mat{I}$; the
$h_{w,i}$ are the wheels' spin momenta relative to the body). Under
zero external torque the inertially resolved
$\mat{R}(\quat{q})\,\vv{H}$ is exactly conserved for any internal wheel
torque profile, because
$\dot{\vv{H}} + \vv{\omega} \times \vv{H} = \vv{0}$: the wheel terms
$\pm\tau \hat{\vv{a}}$ cancel between $\mat{I}\dot{\vv{\omega}}$ and
$\dot{h}_w \hat{\vv{a}}$. This identity is the Phase 4
exit-criterion-7 conservation gate and the reason the bookkeeping
function ships with the model rather than with the tests.

\section{Discretization and implementation}
\label{sec:actuators:implementation}

\begin{enumerate}
  \item \textbf{Module and traceability.} Both actuators live in
    \texttt{cpp/src/models/actuators.cpp}; the source comments echo the
    labels \texttt{eq:actuators:mib} through
    \texttt{eq:actuators:totalmomentum} verbatim.
  \item \textbf{Exact thresholds and rails.} The MIB branch
    \eqref{eq:actuators:mib} returns literal zeros below the threshold,
    and the rail branches of \eqref{eq:actuators:wheelclamp} assign the
    configured limit itself (not an accumulated sum), so refusal,
    saturation, and rail states are bit-exact and reproducible (D-10).
  \item \textbf{Pure functions over plain structs.} Thruster and wheel
    parameters are trivially constructible structs mirroring the FR-13
    vocabulary; the wheel update is
    $(\text{params}, \text{command}, \text{state}, \Delta t) \mapsto
    (\text{delivered}, \text{state}')$ with no globals or I/O (D-2).
    Cluster sums iterate thrusters in configuration order.
  \item \textbf{Determinism.} No libm calls anywhere in the module;
    every operation is an IEEE-754 basic operation in fixed order, so
    the module is bit-portable across platforms under the D-10 flags.
\end{enumerate}

\section{Validation evidence}
\label{sec:actuators:validation}

Table~\ref{tab:actuators:validation} lists the tests that constitute
this chapter's validation evidence. Each test identifier is
machine-checked to exist in the test tree by
\texttt{scripts/lint\_chapters.py}; golden provenance and tolerances
are recorded in \texttt{tests/golden/actuators/manifest.toml}.

\begin{table}[htbp]
  \centering
  \caption{Validation evidence for the actuator models (doctest,
    \texttt{cpp/tests/test\_actuators.cpp}).}
  \label{tab:actuators:validation}
  \begin{tabular}{lp{8.6cm}}
    \toprule
    Test ID & Acceptance basis \\
    \midrule
    \testid{ACT-RCS-MIB} &
      A pulse below the configured MIB delivers exactly zero impulse in
      every component, and the committed $2\times$-MIB pulse matches
      the spec impulse $2 J_{\min}$ to $10^{-12}$ relative (Phase 4
      exit criterion 7); the exact-boundary pulse ($F T = J_{\min}$
      with exactly representable inputs) delivers in full. \\
    \testid{ACT-RCS-COUPLING-GOLDEN} &
      Committed pulse geometries match the independent 60-digit mpmath
      evaluation of \eqref{eq:actuators:mib} and
      \eqref{eq:actuators:rcscoupling} to $10^{-12}$ norm-relative on
      linear and angular impulse, including off-CG and retrograde
      directions and a cluster sum. \\
    \testid{ACT-WHEEL-SATURATION} &
      Torque and momentum clamp exactly at the configured saturations
      per \eqref{eq:actuators:wheelclamp}: over-limit commands deliver
      exactly $\pm\tau_{\max}$, rail landings set the momentum exactly
      to $\pm h_{\max}$, a same-sign command on the rail delivers
      exactly zero, and desaturating commands act normally (Phase 4
      exit criterion 7), all against committed golden vectors. \\
    \testid{ACT-WHEEL-MOMENTUM-CONSERVATION} &
      A three-axis rigid-body slew driven through the wheel model
      (RK4, quaternion attitude, gyroscopic coupling active) conserves
      the inertially resolved total momentum
      \eqref{eq:actuators:totalmomentum} to $10^{-12}$ relative over
      the full maneuver, and the reaction sign
      \eqref{eq:actuators:wheelreaction} moves body and wheel momentum
      in opposite directions (Phase 4 exit criterion 7). \\
    \testid{ACT-DOMAIN} &
      Out-of-domain inputs per Section~\ref{sec:actuators:domain}
      throw \texttt{std::domain\_error}. \\
    \bottomrule
  \end{tabular}
\end{table}

\section{References}
\label{sec:actuators:references}

The momentum-exchange formulation --- wheel momentum states, the
body-reaction sign convention, and the total-momentum decomposition
\eqref{eq:actuators:totalmomentum} --- follows Markley and
Crassidis~\cite{markley2014}. The MIB threshold and the
saturation-rail semantics are configuration-defined device limits
specified in Section~\ref{sec:actuators:domain} of this chapter rather
than external physical results; the impulse-coupling algebra is derived
in Section~\ref{sec:actuators:derivation} from the definition of
force and moment impulses. Full bibliographic details appear in the
bibliography.
